% =============================================================================================
\chapter{Fondements théoriques}\label{chap:theorie}

Ce chapitre présente les concepts, méthodes et algorithmes utilisés, dans l'ordre de la chaîne de décision :
le risque de crédit et les indicateurs financiers, l'apprentissage supervisé et son évaluation, l'interprétabilité,
l'aide multicritère à la décision, la sélection sous budget et les agents fondés sur des modèles de langage.

% ---------------------------------------------------------------------------------------------
\section{Risque de crédit}

\subsection{Défaut et probabilité de défaut}

Un emprunteur fait \emph{défaut} lorsqu'il ne rembourse pas sa dette et que la banque constate une perte. Dans les
données de la SBA, un prêt passé en perte porte le statut \code{CHGOFF} (\textit{charged off}). La \emph{probabilité de
défaut} (PD) est la probabilité qu'un emprunteur fasse défaut sur un horizon donné. Avec la perte en cas de défaut
et l'exposition, elle compose la perte attendue, au cœur de la gestion du risque de crédit et de la réglementation
bâloise~\cite{bcbs2006}.

\begin{definition}[Variable cible]
Pour chaque prêt $i$, on définit $y_i = 1$ si le prêt est passé en perte au plus 60 mois après son premier
versement, et $y_i = 0$ sinon. L'horizon de cinq ans correspond à celui d'un crédit d'investissement.
\end{definition}

\subsection{Scoring de crédit}

Le \emph{scoring} attribue à chaque demande un score de risque à partir de caractéristiques connues au moment de la
demande~\cite{thomas2002}. On attend du score deux qualités distinctes : bien \emph{ordonner} les emprunteurs du plus
sûr au plus risqué (discrimination), et donner des probabilités \emph{fiables} au sens où, parmi les emprunteurs à
qui l'on attribue 5\,\%, environ 5\,\% font défaut (calibration). Dans notre système, la seconde qualité est aussi
importante que la première, car la méthode multicritère et le filtre éliminatoire utilisent la valeur de la
probabilité, pas seulement le rang.

% ---------------------------------------------------------------------------------------------
\section{Indicateurs financiers d'un projet}\label{sec:indicateurs}

\subsection{Valeur actuelle nette}

Un projet nécessite un investissement initial $I_0$ et génère des flux de trésorerie $F_t$ aux années
$t = 1, \dots, T$. Un dirham reçu dans le futur vaut moins qu'un dirham aujourd'hui : on \emph{actualise} les flux au
taux $r$, égal au coût du capital du secteur. La valeur actuelle nette est
\begin{equation}
  \mathrm{VAN} = -I_0 + \sum_{t=1}^{T} \frac{F_t}{(1+r)^t} + \frac{V_T}{(1+r)^T},
\end{equation}
où $V_T$ est la valeur résiduelle du projet après l'horizon explicite. Une VAN positive signifie que le projet crée
de la valeur au-delà de la rémunération exigée par les apporteurs de capitaux. Le taux de rendement interne (TRI) est
le taux $r$ qui annule la VAN.

\subsection{Délai de récupération}

Le délai de récupération est le nombre d'années nécessaires pour que les flux cumulés remboursent
l'investissement :
\begin{equation}
  D = \min \Big\{ \tau : \sum_{t=1}^{\tau} F_t \geq I_0 \Big\},
\end{equation}
interpolé linéairement à l'intérieur de l'année de récupération. C'est une mesure de liquidité et de risque : plus
il est court, plus la banque est vite exposée à un projet qui a déjà remboursé sa mise.

\subsection{Ratio de couverture du service de la dette}

Le DSCR rapporte le flux de trésorerie d'une année au service de la dette de cette année (intérêts et
remboursement du capital) :
\begin{equation}
  \mathrm{DSCR}_t = \frac{F_t}{S_t}, \qquad S = K\,\frac{i}{1-(1+i)^{-n}},
\end{equation}
où $S$ est l'annuité constante d'un prêt de montant $K$ au taux $i$ sur $n$ années. Un DSCR inférieur à 1 signifie
que le projet ne génère pas assez de trésorerie pour payer ses échéances. On retient le minimum sur la durée du
prêt, $\mathrm{DSCR}_{\min} = \min_t \mathrm{DSCR}_t$, qui décrit l'année la plus tendue. Les financements de
projets d'énergie renouvelable exigent typiquement un DSCR de 1,25 à 1,40~\cite{nrel2024}.

\subsection{Levier, garanties et critères extra-financiers}

Le \emph{levier} rapporte la dette aux fonds propres apportés par les propriétaires : plus il est élevé, moins
l'entreprise peut absorber une perte. La \emph{couverture par les garanties} rapporte la valeur des sûretés au
montant prêté ; une couverture de 1 signifie que les garanties valent le prêt. Les orientations de
l'ABE~\cite{eba2020} rappellent que la garantie est un second recours qui ne doit pas fonder la décision, d'où un
poids modéré. Le \emph{score ESG} résume les risques environnementaux et sociaux ; les Principes de
l'Équateur~\cite{equator2020} classent les projets en catégories A (impacts importants), B (limités) et C (minimes).
Enfin, le nombre d'\emph{emplois soutenus par million investi} mesure l'impact économique local.

% ---------------------------------------------------------------------------------------------
\section{Apprentissage supervisé}

On dispose de $n$ exemples $(\bm{x}_i, y_i)$, où $\bm{x}_i \in \mathbb{R}^d$ décrit le prêt et $y_i \in \{0,1\}$
son issue. On cherche une fonction $\hat p(\bm{x}) \approx \Pr(y = 1 \mid \bm{x})$ qui se généralise à des prêts
futurs.

\subsection{Régression logistique}

La régression logistique modélise le logarithme de la cote par une combinaison linéaire des variables :
\begin{equation}
  \log \frac{\hat p(\bm{x})}{1-\hat p(\bm{x})} = \beta_0 + \bm{\beta}^\top \bm{x},
  \qquad \hat p(\bm{x}) = \frac{1}{1+e^{-(\beta_0 + \bm{\beta}^\top \bm{x})}}.
\end{equation}
Les coefficients minimisent la log-vraisemblance négative régularisée :
\begin{equation}
  \min_{\beta_0,\bm{\beta}} \; -\sum_{i=1}^{n} \Big[ y_i \log \hat p(\bm{x}_i) + (1-y_i)\log\big(1-\hat p(\bm{x}_i)\big) \Big]
  + \frac{1}{2C}\,\|\bm{\beta}\|_2^2 .
\end{equation}
Le paramètre $C$ règle la force de la régularisation (plus $C$ est petit, plus les coefficients sont contraints). Les
variables catégorielles sont codées en indicatrices (\textit{one-hot}) et les variables numériques standardisées.
Simple, stable et lisible, c'est le modèle de référence (\textit{baseline}) du scoring de crédit~\cite{thomas2002}.
Sa limite est de ne capter ni non-linéarité ni interaction sans ingénierie manuelle.

\subsection{Modèles additifs généralisés et EBM}

Un modèle additif généralisé (GAM) remplace chaque terme linéaire par une fonction quelconque d'une variable :
\begin{equation}
  g\big(\mathbb{E}[y \mid \bm{x}]\big) = \beta_0 + \sum_{j=1}^{d} f_j(x_j) + \sum_{(j,k) \in \mathcal{I}} f_{jk}(x_j, x_k),
\end{equation}
où $g$ est la fonction logit et $\mathcal{I}$ un petit ensemble de paires de variables en interaction (modèle
GA$^2$M~\cite{lou2013}). L'\emph{Explainable Boosting Machine} (EBM) de la bibliothèque InterpretML~\cite{nori2019}
apprend chaque fonction $f_j$ par \emph{boosting cyclique} : à chaque tour, elle ajuste un arbre très peu profond sur
une seule variable à la fois, avec un faible taux d'apprentissage, en tournant sur toutes les variables. Les
interactions les plus utiles sont ensuite détectées et ajoutées. L'ensemble est moyenné sur plusieurs
rééchantillonnages (\textit{outer bags}).

Le modèle obtenu est aussi précis que de nombreux modèles opaques, tout en restant une \emph{boîte de verre} : la
prédiction est la somme de contributions lisibles, et chaque fonction $f_j$ peut être tracée pour montrer comment le
risque varie avec la variable. C'est la propriété qui motive son choix pour une décision de crédit.

\subsection{Boosting de gradient et LightGBM}

Le boosting de gradient~\cite{friedman2001} construit un modèle additif d'arbres
$F_M(\bm{x}) = \sum_{m=1}^{M} \nu\, h_m(\bm{x})$ : chaque nouvel arbre $h_m$ est ajusté sur le gradient négatif de
la fonction de perte (ici la perte logistique) par rapport à la prédiction courante, et $\nu$ est le taux
d'apprentissage. LightGBM~\cite{ke2017} est une implémentation très efficace : il discrétise les variables en
histogrammes, fait croître les arbres feuille par feuille (la feuille qui réduit le plus la perte d'abord), traite
nativement les variables catégorielles et propose des techniques d'échantillonnage des observations (GOSS) et de
regroupement de variables (EFB). Ses hyperparamètres principaux sont le nombre de feuilles, le taux d'apprentissage
et le nombre d'arbres, ce dernier étant fixé par \emph{arrêt anticipé} : on arrête d'ajouter des arbres quand la
performance sur un jeu de validation cesse de s'améliorer. Très performant sur les données tabulaires, ce modèle
est en revanche opaque : des centaines d'arbres en interaction ne se lisent pas.

% ---------------------------------------------------------------------------------------------
\section{Évaluation d'un classifieur probabiliste}

\subsection{Matrice de confusion, précision, rappel et F1}

Pour un seuil $s$, on prédit un défaut quand $\hat p(\bm{x}) \geq s$. On compte les vrais positifs (VP), faux
positifs (FP), faux négatifs (FN) et vrais négatifs (VN). La précision $\mathrm{VP}/(\mathrm{VP}+\mathrm{FP})$ est la
part des alertes justes, le rappel $\mathrm{VP}/(\mathrm{VP}+\mathrm{FN})$ la part des défauts détectés, et
\begin{equation}
  F_1 = \frac{2\,\text{précision}\times\text{rappel}}{\text{précision}+\text{rappel}} .
\end{equation}
Le F1 dépend du seuil : nous le choisissons sur la validation, là où le F1 est maximal, puis l'appliquons tel quel
au test.

\subsection{Courbes ROC et précision-rappel}

La courbe ROC trace le taux de vrais positifs contre le taux de faux positifs pour tous les seuils. Son aire, l'AUROC,
a une interprétation directe : c'est la probabilité qu'un emprunteur en défaut tiré au hasard reçoive un score plus
élevé qu'un emprunteur sain tiré au hasard ; 0,5 correspond au hasard et 1 à un classement parfait. Quand la classe
positive est rare, l'AUROC peut paraître élevée alors que les alertes sont peu précises : l'aire sous la courbe
précision-rappel (AUPRC) est alors plus informative~\cite{davis2006}. Sa valeur de référence est le taux de défaut :
un modèle au hasard obtient une AUPRC égale à la proportion de positifs.

\subsection{Score de Brier et calibration}

Le score de Brier~\cite{brier1950} mesure l'écart quadratique moyen entre la probabilité prédite et l'issue :
\begin{equation}
  \mathrm{BS} = \frac{1}{n} \sum_{i=1}^{n} \big(\hat p(\bm{x}_i) - y_i\big)^2 .
\end{equation}
Il récompense à la fois la discrimination et la calibration. Un modèle est \emph{calibré} si
$\Pr(y=1 \mid \hat p = p) = p$ ; on le vérifie par une courbe de calibration qui compare, par tranche de probabilité
prédite, la moyenne prédite et la fréquence observée.

La \emph{calibration isotonique}~\cite{zadrozny2002,niculescu2005} corrige les probabilités par une fonction
croissante $m$ ajustée sur un jeu distinct de celui de l'entraînement :
\begin{equation}
  \hat m = \argmin_{m \text{ croissante}} \sum_{i} \big(y_i - m(\hat p_i)\big)^2 ,
\end{equation}
résolu par l'algorithme \textit{pool adjacent violators}. Elle ne modifie pas l'ordre des prédictions (donc
quasiment pas l'AUROC) mais corrige leur niveau. Elle est indispensable lorsque le taux de défaut de la période
d'entraînement diffère de celui de la période d'usage, ce qui est notre cas.

\subsection{Déséquilibre des classes}

Avec 3 à 13\,\% de défauts selon les périodes, la classe positive est minoritaire. Nous n'avons pas rééquilibré
les données (par sur- ou sous-échantillonnage), car cela fausserait les probabilités dont le MCDA a besoin. Nous
avons préféré des métriques adaptées (AUPRC, Brier), un seuil de décision choisi sur la validation et une calibration.

% ---------------------------------------------------------------------------------------------
\section{Bonnes pratiques : séparation des données et fuites}

\subsection{Séparation temporelle}

Les données sont séparées en trois blocs \emph{avant tout prétraitement} : l'entraînement sert à apprendre, la
validation à régler (hyperparamètres, calibration, seuil, choix du modèle) et le test à mesurer la performance
finale, une seule fois. Une banque prédit l'avenir à partir du passé : la séparation suit donc l'année d'approbation
du prêt, et non un tirage au hasard qui mélangerait des prêts de 2008 et de 2015 et donnerait des performances
trop optimistes. Tout ce qui apprend quelque chose des données (médianes d'imputation, échelles, modalités
connues) est ajusté sur l'entraînement uniquement.

\subsection{Fuite de données}

Une fuite survient lorsqu'une variable contient, directement ou indirectement, une information indisponible au
moment de la prédiction~\cite{kaufman2012}. Le modèle paraît alors excellent mais échouera en usage. Les fuites
évidentes (statut du prêt, date et montant de la perte) se suppriment à la source ; les fuites cachées se détectent
par un résultat trop beau pour être vrai, puis par l'examen de la variable suspecte. Nous en avons trouvé une
(section~\ref{sec:fuite}).

% ---------------------------------------------------------------------------------------------
\section{Interprétabilité}

\subsection{Importance par permutation}

L'importance d'une variable $j$ est la baisse de performance observée quand on mélange au hasard ses valeurs, ce qui
casse son lien avec la cible tout en conservant sa distribution~\cite{breiman2001,fisher2019} :
\begin{equation}
  I_j = \mathrm{AUROC}(\hat p, X) - \frac{1}{K}\sum_{k=1}^{K} \mathrm{AUROC}\big(\hat p, X^{\pi_k(j)}\big),
\end{equation}
où $X^{\pi_k(j)}$ est le jeu de données dont la colonne $j$ a été permutée. La méthode s'applique à tout modèle.

\subsection{Explications natives de l'EBM}

L'EBM étant additif, l'importance d'un terme est la valeur absolue moyenne de sa contribution
$\frac{1}{n}\sum_i |f_j(x_{ij})|$, et l'effet d'une variable se lit directement sur la courbe $f_j$. Pour un prêt
donné, la prédiction se décompose exactement en contributions : on peut nommer les facteurs qui augmentent son risque.

\subsection{Valeurs SHAP}

Pour un modèle opaque, les valeurs SHAP~\cite{lundberg2017} attribuent à chaque variable une part de l'écart entre la
prédiction et la prédiction moyenne, selon la valeur de Shapley de la théorie des jeux coopératifs :
\begin{equation}
  \phi_j(\bm{x}) = \sum_{S \subseteq \{1,\dots,d\}\setminus\{j\}}
  \frac{|S|!\,(d-|S|-1)!}{d!}\Big[ v_{\bm{x}}(S\cup\{j\}) - v_{\bm{x}}(S) \Big],
\end{equation}
où $v_{\bm{x}}(S)$ est la prédiction attendue lorsque seules les variables de $S$ sont connues. Ces valeurs sont les
seules à vérifier simultanément l'efficacité (elles somment à l'écart total), la symétrie et la nullité. Pour les
modèles d'arbres, l'algorithme \textit{TreeSHAP} les calcule exactement en temps polynomial.

% ---------------------------------------------------------------------------------------------
\section{Aide multicritère à la décision}

Une alternative $a$ (un projet) est évaluée sur $m$ critères $g_1(a), \dots, g_m(a)$, chacun à maximiser ou à
minimiser, avec des poids $w_j$. Les méthodes \emph{compensatoires}, comme la somme pondérée, permettent qu'une
mauvaise note soit rachetée par de bonnes notes ailleurs. Les méthodes de \emph{surclassement} de l'école
européenne~\cite{roy1991,figueira2005} limitent cette compensation : une très mauvaise performance sur un critère
peut opposer un \emph{veto}. Pour un dossier de crédit, cette propriété est essentielle : un DSCR très faible ne
devrait pas être effacé par une VAN élevée.

\subsection{Pondération par AHP}

L'\textit{Analytic Hierarchy Process}~\cite{saaty1980} construit les poids à partir de comparaisons deux à deux. Pour
chaque paire de critères $(i,j)$, le décideur indique l'importance relative $a_{ij}$ sur l'échelle de Saaty (1 :
égale importance, 3 : modérée, 5 : forte, 7 : très forte, 9 : extrême, valeurs paires intermédiaires), avec
$a_{ji} = 1/a_{ij}$. Huit critères demandent $8\times7/2 = 28$ jugements. Le vecteur des poids est le vecteur propre
principal de la matrice $A = (a_{ij})$ :
\begin{equation}
  A\,\bm{w} = \lambda_{\max}\,\bm{w}, \qquad \sum_j w_j = 1 .
\end{equation}
Si les jugements étaient parfaitement cohérents ($a_{ik} = a_{ij}a_{jk}$), on aurait $\lambda_{\max} = m$. L'indice
de cohérence $\mathrm{CI} = (\lambda_{\max}-m)/(m-1)$ est rapporté à un indice aléatoire $\mathrm{RI}$ ($1{,}41$ pour
$m = 8$) : le ratio $\mathrm{CR} = \mathrm{CI}/\mathrm{RI}$ doit rester inférieur à 0,10.

\subsection{ELECTRE III}

ELECTRE III~\cite{roy1978} compare les alternatives deux à deux. Chaque critère $j$ reçoit trois seuils : un écart
inférieur au seuil d'indifférence $q_j$ est négligeable, un écart supérieur au seuil de préférence $p_j$ est net, et
un écart supérieur au seuil de veto $v_j$ est rédhibitoire. Pour des critères à maximiser, on note
$\Delta_j = g_j(b) - g_j(a)$ l'avantage de $b$ sur $a$.

\paragraph{Concordance.} Le critère $j$ soutient l'affirmation « $a$ surclasse $b$ » d'autant plus que $b$ ne fait
pas nettement mieux :
\begin{equation}
  c_j(a,b) =
  \begin{cases}
    1 & \text{si } \Delta_j \leq q_j,\\[2pt]
    \dfrac{p_j - \Delta_j}{p_j - q_j} & \text{si } q_j < \Delta_j < p_j,\\[6pt]
    0 & \text{si } \Delta_j \geq p_j,
  \end{cases}
  \qquad
  C(a,b) = \frac{\sum_j w_j\, c_j(a,b)}{\sum_j w_j}.
\end{equation}

\paragraph{Discordance.} Le critère $j$ s'oppose d'autant plus que $b$ fait beaucoup mieux :
\begin{equation}
  d_j(a,b) =
  \begin{cases}
    0 & \text{si } \Delta_j \leq p_j,\\[2pt]
    \dfrac{\Delta_j - p_j}{v_j - p_j} & \text{si } p_j < \Delta_j < v_j,\\[6pt]
    1 & \text{si } \Delta_j \geq v_j .
  \end{cases}
\end{equation}

\paragraph{Crédibilité.} La concordance est atténuée par les critères dont la discordance la dépasse :
\begin{equation}
  \sigma(a,b) = C(a,b) \prod_{j \,:\, d_j(a,b) > C(a,b)} \frac{1 - d_j(a,b)}{1 - C(a,b)} .
\end{equation}
Un veto ($d_j = 1$) annule la crédibilité, quelle que soit la concordance : c'est la non-compensation.

\paragraph{Distillation.} Le classement s'obtient en exploitant la matrice $\sigma$ par deux distillations.
À chaque étape, on fixe un niveau $\lambda_k$ juste en dessous de la plus forte crédibilité restante, avec un seuil
de discrimination $s(\lambda) = 0{,}30 - 0{,}15\,\lambda$, et l'on considère que $a$ surclasse $b$ si
$\sigma(a,b) > \lambda_k$ et $\sigma(a,b) - \sigma(b,a) > s(\sigma(a,b))$. La qualification d'une alternative est
le nombre d'alternatives qu'elle surclasse moins le nombre de celles qui la surclassent. La distillation
\emph{descendante} extrait successivement les meilleures, l'\emph{ascendante} les moins bonnes. Les deux préordres
peuvent différer : une alternative bien placée par l'une et mal placée par l'autre est \emph{incomparable} à une
partie des autres, information propre aux méthodes de surclassement.

\subsection{ELECTRE TRI}

ELECTRE TRI~\cite{yu1992} affecte chaque alternative à une catégorie ordonnée, au lieu de la classer. Les catégories
sont séparées par des \emph{profils limites} : $b_1$ entre Refuser et Conditionner, $b_2$ entre Conditionner et
Financer. On calcule la crédibilité $\sigma(a, b_h)$ et $\sigma(b_h, a)$ comme dans ELECTRE III, puis on la
compare à un niveau de \emph{coupure} $\lambda \in [0{,}5 ; 1]$ : $a$ surclasse $b_h$ si $\sigma(a,b_h) \geq \lambda$.

\begin{itemize}
  \item \textbf{Règle pessimiste} (prudente) : on compare $a$ aux profils du plus exigeant au moins exigeant, et on
        l'affecte à la catégorie juste au-dessus du premier profil qu'il surclasse.
  \item \textbf{Règle optimiste} : on compare $a$ aux profils du moins exigeant au plus exigeant, et on l'affecte à la
        catégorie juste en dessous du premier profil qui lui est strictement préféré.
\end{itemize}
La règle pessimiste n'admet un projet dans une catégorie que s'il en franchit clairement le seuil d'entrée ; elle
correspond à la prudence bancaire.

\begin{figure}[htbp]
\centering
\begin{tikzpicture}[x=1.3cm,y=0.62cm,font=\small]
  \foreach \j/\lab in {1/PD,2/DSCR,3/VAN,4/Délai,5/Levier,6/Garanties,7/ESG,8/Emplois}{
    \draw[eccgris!60] (\j,0) -- (\j,6);
    \node[below,text=eccfonce] at (\j,0) {\lab};}
  \fill[rouge!12] (0.6,0) rectangle (8.4,2);
  \fill[ambre!14] (0.6,2) rectangle (8.4,4);
  \fill[eccvert!12] (0.6,4) rectangle (8.4,6);
  \foreach \j in {1,...,8}{\draw[eccgris!60] (\j,0) -- (\j,6);}
  \draw[rouge,thick] (1,2.1) -- (2,1.8) -- (3,2.2) -- (4,1.9) -- (5,2.1) -- (6,1.8) -- (7,2.0) -- (8,2.2);
  \draw[eccvert,thick] (1,4.0) -- (2,4.2) -- (3,3.8) -- (4,4.1) -- (5,3.9) -- (6,4.2) -- (7,4.0) -- (8,3.9);
  \draw[eccfonce,thick,dashed,mark=*,mark options={solid,fill=eccfonce}] plot coordinates
    {(1,4.8) (2,3.1) (3,1.2) (4,1.5) (5,4.6) (6,3.3) (7,5.2) (8,2.9)};
  \node[right,text=rouge] at (8.45,2.1) {profil $b_1$};
  \node[right,text=eccvert] at (8.45,4.0) {profil $b_2$};
  \node[right,text=eccfonce] at (8.45,5.4) {projet $a$};
  \node[left,text=eccvert] at (0.55,5) {Financer};
  \node[left,text=ambre] at (0.55,3) {Conditionner};
  \node[left,text=rouge] at (0.55,1) {Refuser};
\end{tikzpicture}
\caption{Principe d'ELECTRE TRI : chaque critère (axe vertical, \og meilleur \fg{} vers le haut) est découpé par les
profils $b_1$ et $b_2$. Le projet $a$ est comparé globalement à chaque profil par la crédibilité $\sigma$, et non
critère par critère.}
\label{fig:electretri}
\end{figure}

\subsection{PROMETHEE II}

PROMETHEE~\cite{brans1985} associe à chaque critère une fonction de préférence $P_j(a,b) \in [0,1]$ croissante avec
l'écart $\delta_j = g_j(a) - g_j(b)$. Nous utilisons la fonction linéaire avec zone d'indifférence :
$P_j = 0$ si $\delta_j \leq q_j$, $P_j = 1$ si $\delta_j \geq p_j$ et
$P_j = (\delta_j - q_j)/(p_j - q_j)$ entre les deux. L'indice de préférence agrégé et les flux sont
\begin{equation}
  \pi(a,b) = \sum_j w_j P_j(a,b), \qquad
  \phi^+(a) = \frac{1}{n-1}\sum_{b\neq a}\pi(a,b), \qquad
  \phi^-(a) = \frac{1}{n-1}\sum_{b\neq a}\pi(b,a),
\end{equation}
et PROMETHEE II classe les alternatives par flux net décroissant $\phi(a) = \phi^+(a) - \phi^-(a)$. La méthode ne
comporte pas de veto : nous l'utilisons pour vérifier que le classement ne dépend pas de la méthode choisie.

\subsection{Comparer deux classements}

La concordance entre deux classements de $n$ alternatives se mesure par le coefficient de Spearman,
$\rho = 1 - \frac{6\sum_i d_i^2}{n(n^2-1)}$ où $d_i$ est l'écart de rang de l'alternative $i$, et par le $\tau$ de
Kendall, qui compte la proportion de paires classées dans le même ordre moins la proportion de paires inversées.
Les deux valent 1 pour des classements identiques et 0 en l'absence de lien.

% ---------------------------------------------------------------------------------------------
\section{Sélection sous budget}\label{sec:theorie-budget}

Choisir quels projets financer sous un budget $B$ est un \emph{problème du sac à dos}~\cite{martello1990}. Chaque
projet $i$ demande un montant $a_i$ et apporte une valeur $v_i$ ; la variable $x_i \in \{0,1\}$ indique s'il est
retenu :
\begin{equation}
  \max_{\bm{x}\in\{0,1\}^n} \sum_i v_i\,x_i
  \quad\text{s.c.}\quad \sum_i a_i\,x_i \leq B,
  \quad \sum_{i \in s} a_i\,x_i \leq \alpha B \;\; \forall s \in \mathcal{S},
  \quad x_i = 0 \text{ si } a_i \notin [a_{\min}, a_{\max}],
\end{equation}
où $\mathcal{S}$ est l'ensemble des secteurs et $\alpha$ la part maximale d'un secteur. Ce programme linéaire en
nombres entiers est NP-difficile en général, mais se résout exactement pour quelques centaines de projets par
séparation et évaluation (\textit{branch and bound}) ; nous utilisons le solveur HiGHS~\cite{huangfu2018}.

Le choix des valeurs $v_i$ détermine la politique de sélection. Si l'on veut qu'un projet mieux classé passe
\emph{toujours} avant tous les projets moins bien classés, on donne des poids lexicographiques
($v_i > \sum_{k \text{ moins bien classé}} v_k$). Le problème se résout alors exactement de façon gloutonne :
on parcourt les projets dans l'ordre du classement et l'on retient chacun s'il respecte encore les contraintes.

% ---------------------------------------------------------------------------------------------
\section{Modèles de langage et agents}

\subsection{Grands modèles de langage}

Un grand modèle de langage (LLM) est un réseau de neurones de type \textit{transformer}, entraîné sur de très
grands corpus à prédire le mot suivant, puis ajusté pour suivre des instructions. Il rédige, résume et explique avec
aisance. Il peut aussi produire des affirmations fausses formulées avec assurance (\emph{hallucinations}),
y compris des chiffres inventés~\cite{ji2023}. Les LLM calculent mal et de façon non reproductible : ils ne doivent
pas produire les nombres d'une décision de crédit.

\subsection{Agents et appel d'outils}

Un \emph{agent} est un programme qui utilise un LLM pour accomplir une tâche en plusieurs étapes, en raisonnant et
en agissant sur son environnement~\cite{yao2023}. L'\emph{appel d'outils} (\textit{function calling}) permet au LLM
de demander l'exécution de fonctions décrites par un schéma (nom, rôle, paramètres), puis de rédiger à partir de
leurs résultats~\cite{schick2023}. Le calcul reste ainsi dans du code déterministe et vérifiable, et le LLM se limite
au langage : comprendre une question, choisir les bons outils, reformuler les résultats.

\subsection{Garde-fous}

Trois garde-fous complètent ce schéma. Le LLM ne reçoit que des données déjà calculées ou n'y accède que par des
outils. Un vérificateur extrait les nombres de sa réponse et contrôle qu'ils figurent dans les données consultées.
Enfin, la décision reste humaine : l'agent conseille, le comité décide.
